\section{Data and Frequent Losers Classification}
\label{sec:data}

\subsection{Data}
\label{sec:data_source}

The data come from the BEC platform described in
Section~\ref{sec:bec_platform}. We observe 4.5 million tender-items
and 40 million individual bids from 41,000 firms across 1,308
purchasing units over 2009--2019. Each observation records the item
description, product code, modality (convite or preg\~ao), procuring
unit, date, all bid amounts, and the winning price. The data include
losing bids---a feature essential for identifying participation
patterns.

\subsection{Frequent Losers Classification}
\label{sec:fl_definition}

We identify frequent losers (FL) in two steps, following
\citet{genicolomartins2026screening}:

\begin{enumerate}
  \item \textbf{Always-losers}: firms with a zero win rate across all
    2009--2019 tenders. Of 41,000 BEC participants, 16,843 ($\sim$41\%)
    never win. The high always-loser rate reflects small firms bidding
    beyond capacity, BEC's low participation cost, and---in
    convite---PBU officials inviting firms to satisfy the three-bidder
    rule.

  \item \textbf{IQR threshold}: among always-losers, we flag those
    with abnormally high participation using a
    median~$+$~1.5~$\times$~IQR threshold on the tender-count
    distribution.%
    \footnote{We use the median rather than $Q_3$ because the
    distribution is heavily right-skewed. See
    \citet{genicolomartins2026screening} for robustness to
    alternative multipliers (1.0--3.0$\times$IQR).}
    The resulting cutoff is 14 tenders. The economic justification
    is the negative-profit argument in Section~\ref{sec:rationality}:
    at 14+ tenders with zero wins, expected profit is strictly
    negative for any positive bidding cost, so persistent
    participation requires a non-competitive explanation.
    This yields \textbf{2,735 FL firms}.
\end{enumerate}

\noindent The treatment variable $\text{losers}_{igt}$ equals one if
tender-item $i$ in item group $g$ at time $t$ has at least one FL
firm among its participants.

We note three properties of the classification. \emph{First}, it is
static: FL status is defined over the full 2009--2019 window.
Splitting the sample at 2013, only 10\% of first-half FL firms
reappear as FL in the second half, suggesting that the identities
rotate even as the pattern persists. \emph{Second}, it is validated
against external data: FL firms co-participate with CADE-convicted
cartelists at 3.5$\times$ the baseline rate ($p < 0.001$), and the
screen achieves AUC~$= 0.94$ against enforcement outcomes
\citep{genicolomartins2026screening}. \emph{Third}, it is
intentionally coarse. FL is a triage tool that flags suspicious
environments, not a classifier that labels individual firms as
cartel members.

\subsection{Sample Construction}
\label{sec:sample}

The regression sample retains convite and preg\~ao tenders with a
recorded winner, restricted to item types with at least one
FL-present tender. The restriction ensures that controls come from
product markets where FL firms could plausibly participate; the
unrestricted sample yields a nearly identical coefficient.%
\footnote{See \citet{genicolomartins2026screening}, Table~B.3.}

Table~\ref{tab:sample_by_modality} reports the sample composition.
Two features merit attention. \emph{First}, convite accounts for
67\% of tender-items but only 33\% of procurement value---it is the
high-volume, low-value modality where the minimum-bidder rule
applies. \emph{Second}, FL presence rates differ across modalities:
4.8\% overall, with higher incidence in convite tenders where the
rule creates mechanical demand for cover bidders.

\begin{table}[htbp]
\centering
\caption{Sample composition by procurement modality}
\label{tab:sample_by_modality}
\begin{adjustbox}{max width=\textwidth}
\begin{threeparttable}
\small
\begin{tabular}{lccc}
\toprule
 & Preg\~ao & Convite & Total \\
\midrule
Tender-items & 546,549 & 1,107,852 & 1,654,401 \\
Item types & 12,431 & 14,672 & 18,783 \\
Purchasing units (PBUs) & 892 & 1,104 & 1,308 \\
FL-present tender-items (\%) & 3.9 & 5.2 & 4.8 \\
\midrule
\multicolumn{4}{l}{\textit{Among FL-present convite tenders:}} \\
\quad $n_{\text{genuine}} < 3$ (rule binding) & --- & 39.2\% & --- \\
\quad $n_{\text{genuine}} = 0$ (pure cover bidding) & --- & 23.5\% & --- \\
\midrule
Mean log price & 6.82 & 5.47 & 5.91 \\
Mean number of bidders & 4.21 & 3.48 & 3.72 \\
\bottomrule
\end{tabular}
\begin{tablenotes}
\small
\item \textit{Notes:} Sample restricted to tender-items with a
recorded winner in item types with at least one FL-present tender
(2009--2019). $n_{\text{genuine}}$ counts non-FL bidders.
``Rule binding'' means the minimum-bidder threshold ($\geq 3$
participants) could not be met without FL firms. ``Pure cover
bidding'' means all participants are FL firms.
\end{tablenotes}
\end{threeparttable}
\end{adjustbox}
\end{table}

The bottom panel of Table~\ref{tab:sample_by_modality} reports the
two statistics central to the paper's argument. Among FL-present
convite tenders, 39.2\% have fewer than three genuine
competitors---the minimum-bidder rule binds, and the cartel must
deploy cover firms to clear the threshold. In 23.5\% of cases,
$n_{\text{genuine}} = 0$: every participant is a cover firm, and
the rule achieved zero genuine competition at positive
administrative cost. These figures describe the corner solution of
the model in Section~\ref{sec:corner_interior}.
